%==============================================================================
% PAPER A — SOFTWARE / FRAMEWORK PAPER
% CosimGym: a declarative, model-agnostic framework for RL on HELICS co-simulation
% Target: SoftwareX (short) or SIMPAT (full).
% STATUS: first prose draft. Architecture/method sections are real draft text.
% Results/validation contain real numbers ONLY after experiments -> marked [TODO].
% Citations are \cite{...} placeholders; fill cosimgym.bib (see PAPER_PLAN.md §8).
%==============================================================================
\documentclass[preprint,12pt]{elsarticle}
\usepackage{graphicx,booktabs,amsmath,listings,xcolor,hyperref}

\journal{SoftwareX} % or Simulation Modelling Practice and Theory

\begin{document}
\begin{frontmatter}

\title{CosimGym: A Declarative, Model-Agnostic Framework for\\
Reinforcement Learning on HELICS Co-Simulations}

\author[polito]{Pietro Rando Mazzarino\corref{cor}}
\ead{pietro.randomazzarino@polito.it}
\author[polito]{Co-authors}
\address[polito]{EC-Lab, Politecnico di Torino, Turin, Italy}
\cortext[cor]{Corresponding author.}

\begin{abstract}
The decarbonization of energy systems is driving them toward decentralized,
renewable-rich, sector-coupled architectures whose system-level behavior can only
be captured by coupling heterogeneous models across domains and time scales.
Co-simulation has become the established method for this task, while reinforcement
learning (RL) has emerged as a promising data-driven paradigm for adaptive control
of such cyber-physical energy systems. Bridging the two, however, is repeatedly
re-implemented by hand: co-simulation middleware speaks federated
publish/subscribe and federated time, whereas RL tooling speaks the
\texttt{reset()}/\texttt{step()} contract of Gymnasium. Existing RL testbeds are
either locked to a single simulator and domain, or expose a lightweight,
synchronous fixed-frequency environment, limiting realism and reproducibility. We
present \textbf{CosimGym}, a framework that compiles a single declarative YAML
scenario into a running HELICS co-simulation that is simultaneously a standard
Gymnasium environment. Its contributions are a declarative scenario compiler; a
catalog-driven, multi-formalism model abstraction unifying native Python models,
FMI~2.0/3.0 Functional Mock-up Units, EnergyPlus models, data feeders, and RL
agents; native multi-rate, causality-aware, hierarchical multi-federation
orchestration; and a principled treatment of episodic reset over non-resettable
co-simulations. We demonstrate the framework on building thermal control and
PV--battery dispatch with DQN and SAC agents.
% TODO: append one quantitative sentence after experiments.
The framework is open source. % TODO: license + Zenodo DOI.
\end{abstract}

% NOTE: elsarticle has NO \begin{highlights} environment (Elsevier CAS class only).
% For elsarticle, Elsevier wants Highlights as a SEPARATE file at submission.
% Highlights (3-5, <=85 chars each):
%  - Declarative YAML compiled into a HELICS co-simulation that is a Gymnasium env
%  - Model-agnostic catalog unifies Python, FMI 2/3, EnergyPlus, data and RL agents
%  - Native multi-rate, causality-aware, hierarchical multi-federation orchestration
%  - Episodic reset strategies for non-resettable co-simulations (full/soft/rolling/random)
%  - Observation/action spaces auto-derived from physical I/O bounds; telemetry built in

\begin{keyword}
co-simulation \sep reinforcement learning \sep HELICS \sep FMI \sep Gymnasium \sep
cyber-physical energy systems \sep multi-energy systems \sep digital twin
\end{keyword}
\end{frontmatter}

%------------------------------------------------------------------------------
% Optional SoftwareX metadata table (required by that journal). Fill at submission:
% version, repo URL, license, Python 3.12, deps (HELICS, fmpy, Redis, InfluxDB,
% SB3, torch), Zenodo DOI, support email.
%------------------------------------------------------------------------------

%==============================================================================
\section{Motivation and significance}\label{sec:intro}
%==============================================================================
The transition toward decarbonized energy systems is reshaping them into
decentralized, renewable-dominated, and increasingly sector-coupled
infrastructures. Photovoltaic generation, electrified heating, storage, electric
mobility, and active prosumers couple the electrical grid to buildings and to
thermal and gas networks, and these subsystems are most faithfully described by
specialized tools operating at very different time scales---from sub-second
power-electronic dynamics to hourly weather and market signals. Monolithic
simulators cannot represent such cross-domain interactions, and co-simulation has
consequently become the standard methodology for system-level analysis of
cyber-physical energy systems (CPES)~\cite{helics,mosaik}. By coupling models,
tools, and time scales through a common middleware, co-simulation exposes emergent
behavior that is invisible to any single simulator.

In parallel, reinforcement learning (RL) has become a prominent data-driven
paradigm for adaptive control in energy systems, with applications in building
climate control, demand response, energy communities, and grid
operation~\cite{rl_energy_review}. Yet evaluating RL controllers in realistic,
heterogeneous, multi-domain settings remains difficult. The two ecosystems were
designed under different assumptions: co-simulation middleware such as
HELICS~\cite{helics} organizes computation as a federation of processes that
exchange values through publish/subscribe and advance a shared federated clock,
whereas RL libraries such as Gymnasium~\cite{gymnasium},
Stable-Baselines3~\cite{sb3}, and RLlib~\cite{rllib} expect a synchronous,
blocking environment exposing \texttt{reset()} and \texttt{step(action)}. Bridging
this impedance mismatch is, in current practice, re-implemented for every study,
which limits reproducibility, interoperability, and transfer of results across
system configurations.

Existing RL environments only partially close the gap. Building-oriented testbeds
such as Sinergym~\cite{sinergym}, BOPTEST~\cite{boptest}, Energym~\cite{energym},
and CityLearn~\cite{citylearn} are powerful but are tied to a specific simulator
or domain and to largely fixed time stepping. Power-systems environments such as
PowerGridworld~\cite{powergridworld} provide multi-agent Gym interfaces but, by
their authors' own account, rely on synchronous fixed-frequency time stepping and
a limited communication model. Studies that connect HELICS to RL directly, for
example for resilient microgrid control~\cite{frl_microgrid}, build a bespoke
wrapper per experiment that is hard to reuse or reproduce. What is missing is a
\emph{declarative} way to compose arbitrary, multi-rate, multi-formalism physics
behind a standard RL interface---a generalized middleware positioned between pure
co-simulation engines and AI environments, acting as a universal translator
between established engineering tools (EnergyPlus, Modelica/FMU) and modern RL
libraries.

This paper introduces CosimGym, which fills that role. A single declarative YAML
scenario is compiled into a running HELICS federation that is, at the same time, a
standard \texttt{gymnasium.Env}. Its main contributions are:
\begin{itemize}
  \item a \textbf{declarative scenario compiler} that turns YAML into a strictly
        typed configuration and provisions the brokers, processes, timing, and
        data flows of a complete co-simulation without orchestration code;
  \item a \textbf{catalog-driven, multi-formalism model abstraction} that unifies
        native Python models, FMI~2.0/3.0 FMUs, EnergyPlus models, data feeders,
        and RL agents behind one interface;
  \item \textbf{native multi-rate, causality-aware, hierarchical
        multi-federation} orchestration;
  \item a principled treatment of \textbf{episodic reset over non-resettable
        co-simulations}; and
  \item RL \textbf{observation and action spaces auto-derived from physical I/O
        bounds}, with built-in time-series telemetry and dashboards.
\end{itemize}

CosimGym targets researchers and engineers who need to train or benchmark control
policies against realistic physics without rebuilding simulation infrastructure.
The model catalog and scenario compiler are reusable assets that lower the barrier
to reproducible RL-on-physics research, and the design is forward-looking: the same
substrate extends naturally toward multi-agent and multi-timescale learning and
toward sim-to-real deployment (Section~\ref{sec:future}). The remainder of the
paper reviews background and related work (Section~\ref{sec:background}), details
the architecture (Section~\ref{sec:arch}), formalizes the reset problem
(Section~\ref{sec:reset}), reports illustrative validation
(Section~\ref{sec:examples}), discusses impact (Section~\ref{sec:impact}), and
concludes with limitations and future potential
(Sections~\ref{sec:future}--\ref{sec:conclusion}).

%==============================================================================
\section{Background and related work}\label{sec:background}
%==============================================================================
\paragraph{Co-simulation and model-exchange standards}
HELICS (Hierarchical Engine for Large-scale Infrastructure
Co-Simulation)~\cite{helics} couples independent simulators, called federates,
through brokers that mediate value and message exchange and coordinate a shared
notion of federated time; federates may run at different rates and on different
hosts, and brokers can be arranged hierarchically for scale. Mosaik~\cite{mosaik}
addresses a similar need with a Python-centric scheduler. Orthogonally, the
Functional Mock-up Interface (FMI)~\cite{fmi} standardizes the packaging of
dynamic models as Functional Mock-up Units (FMUs), portable containers supported
by more than one hundred tools and routinely used as digital-twin building
blocks. These technologies provide interoperability and time coordination but
deliberately offer no notion of episodes, rewards, or agent--environment
interaction.

\paragraph{RL environments for energy systems}
Gymnasium~\cite{gymnasium} defines the de facto agent--environment contract, and
libraries such as SB3~\cite{sb3} and RLlib~\cite{rllib} consume it. Within energy,
Sinergym~\cite{sinergym} wraps EnergyPlus, BOPTEST~\cite{boptest} exposes Modelica
emulators for building control, Energym~\cite{energym} offers an FMU-based model
library, and CityLearn~\cite{citylearn} targets district-level, multi-agent
demand response at fixed hourly resolution. These environments are valuable but
each is bound to a particular simulator or domain and to a largely fixed temporal
structure.

\paragraph{RL on co-simulation}
A smaller body of work places RL directly on co-simulation.
PowerGridworld~\cite{powergridworld} composes component Gym environments with an
OpenDSS power-flow solver and supports multi-agent training, but its authors
identify synchronous fixed-frequency time stepping and a limited communication
model as limitations. HELICS-based studies couple detailed grid simulators with
RL for specific problems such as resilient microgrid control~\cite{frl_microgrid},
and CommonPower~\cite{commonpower} emphasizes safe, model-based smart-grid
control. Across this landscape, RL environments tend to be domain-specific and
monolithic, while co-simulation tools lack native AI integration. CosimGym is
designed as the generalized, configurable middleware between the two.
Table~\ref{tab:compare} summarizes the comparison.

\begin{table}[t]\centering
\caption{Positioning of CosimGym against representative frameworks.}
\label{tab:compare}
\small
\begin{tabular}{@{}lccccc@{}}
\toprule
Framework & Physics backend & Multi-rate & Multi-domain & MARL & Declarative \\
\midrule
Sinergym       & EnergyPlus           & no  & no  & ltd & partial \\
BOPTEST        & Modelica emulators   & no  & no  & no  & no \\
CityLearn      & built-in data models & no  & district & yes & schema \\
Energym        & Modelica/EnergyPlus  & ltd & no  & no  & no \\
PowerGridworld & OpenDSS + components & no  & no  & yes & partial \\
\textbf{CosimGym} & \textbf{Python/FMI/EnergyPlus} & \textbf{yes} & \textbf{yes} & roadmap & \textbf{yes} \\
\bottomrule
\end{tabular}
\end{table}

%==============================================================================
\section{Architecture and generalized design principles}\label{sec:arch}
%==============================================================================
CosimGym enforces a strict separation of concerns among \emph{configuration},
\emph{execution orchestration}, and \emph{physical modeling}. This separation is
what makes the framework domain-agnostic: the same execution machinery runs a
two-node toy model and a multi-domain federation, and the same RL wrapper trains an
agent on either. Figure~\ref{fig:arch} gives the end-to-end view.

\begin{figure}[t]\centering
% TODO: architecture diagram.
% User YAML -> ScenarioManager (+ Redis state) -> HELICS broker hierarchy ->
% heterogeneous base federates (FMU / Python / data) <-> universal RL federate ->
% telemetry (InfluxDB) + dashboard (Streamlit).
\fbox{\parbox{0.9\linewidth}{\centering\vspace{2.5cm}
\textbf{[Figure 1: architecture / data flow]}\vspace{2.5cm}}}
\caption{From a declarative YAML scenario to a running co-simulation that is also a
Gymnasium environment. The compiler provisions brokers and processes; the model
catalog (synchronized via Redis) supplies model code and metadata; the universal RL
federate exposes the federation as a single environment; telemetry is logged
centrally.}
\label{fig:arch}
\end{figure}

\subsection{Declarative orchestration layer}
The foundation of the framework is that the user describes \emph{what} to simulate,
never \emph{how} to orchestrate it. A YAML scenario declares the federation
topology, each federate's execution rate and publish/subscribe connections, the
model to instantiate, and, for learning scenarios, the RL task. The
\texttt{ScenarioManager} parses this description into strictly typed configuration
objects (validated with Pydantic), provisions the required HELICS broker
hierarchy, and launches each federate as an independent operating-system process
through a common launcher. The full configuration is serialized to a Redis store so
that every process can retrieve its own parameters at startup. Because scaling the
system changes only the declarative description, moving from a single-federation
example to a multi-domain, multi-federation study requires no change to the
execution logic. Listing~\ref{lst:yaml} shows the structure of a minimal scenario.

\begin{lstlisting}[language=,caption={Skeleton of a scenario (abridged).},label=lst:yaml,basicstyle=\ttfamily\footnotesize,frame=single]
name: building_hp_rl
start_time: "2024-01-01T00:00:00"
end_time:   "2024-01-02T00:00:00"
reinforcement_learning_config:
  agent: {model_name: rl_simple_SACsb3, env: {observations: [...], actions: [...]}}
  training: {mode: online, episode_length: 2880, reset_mode: rolling}
federations:
  federation_1:
    federate_configs:
      weather:  {type: base, timing_configs: {real_period: 60}, ...}
      heatpump: {type: base, timing_configs: {real_period: 60}, ...}
      building: {type: base, timing_configs: {real_period: 60}, ...}
\end{lstlisting}

\subsection{Extensible model catalog and state management}
Model code is decoupled from scenarios through a central catalog. Each catalog
entry maps a model name to an implementing class and module together with a typed
specification of its inputs, outputs, and parameters; every variable carries a
type, physical unit, optional bounds, and descriptive tags. At startup the catalog
is loaded into Redis, and at runtime a federate resolves its model metadata,
imports the corresponding class, or unpacks an FMU---loaded from local storage or
pulled from an object store. All models implement the same minimal interface
(\texttt{initialize}, \texttt{step}, \texttt{finalize}), so physical models,
controllers, data readers, and RL agents are treated uniformly. Adding a model
requires only a new class and one catalog entry, which encourages a shared,
reusable library of co-simulation components. Table~\ref{tab:models} lists the
built-in models, which already span several domains and three model formalisms
(native Python, FMI~2.0/3.0, and EnergyPlus exported as an FMU).

\begin{table}[t]\centering
\caption{Representative built-in catalog models.}
\label{tab:models}
\small
\begin{tabular}{@{}lll@{}}
\toprule
Model & Domain & Formalism \\
\midrule
1R1C building, heat pump, PID & building energy & native Python \\
PV, battery, rule-based BEMS   & generation/storage & native Python \\
weather / CSV readers          & data & native Python \\
spring--mass--damper           & generic dynamics & native Python \\
EnergyPlus zone model          & building energy & FMI 2.0 (FMU) \\
type/causality feedthrough     & testing & FMI 3.0 (FMU) \\
DQN, SAC agents                & control & RL agent \\
\bottomrule
\end{tabular}
\end{table}

\subsection{Timing and causality engine}
Realistic CPES studies require federates to advance at genuinely different rates.
Each federate declares a real-world period; the framework normalizes these to
integer HELICS ticks against the smallest period, so a one-minute controller and a
one-hour weather feed coexist without manual bookkeeping. Because federates that
exchange values within the same tick can form algebraic loops, CosimGym lets a
subscription be marked \texttt{same\_step} or \texttt{next\_step}: the former is
consumed immediately, the latter is staged and applied on the following tick,
acting as a one-step delay. Small time offsets can also be applied to break
same-tick cycles, and these offsets can be computed automatically by a topological
sort of the dependency graph; declared causality cycles are validated before
execution. For scenarios with more than one federation, a hierarchy broker is
inserted automatically and ports are assigned dynamically. This timing model is the
co-simulation depth that lightweight, synchronous fixed-frequency
environments~\cite{powergridworld} lack.

\subsection{The universal RL federate wrapper}
A learning scenario adds one further federate that turns the entire federation into
a single \texttt{gymnasium.Env}. The wrapper resolves the impedance mismatch in
three ways. First, observation and action spaces are constructed dynamically from
the scenario mapping and the catalog bounds rather than hard-coded: a bounded
floating-point input becomes a \texttt{Box}, an integer or binned input becomes a
\texttt{Discrete} space, and units and ranges are taken from the model metadata.
Second, actuation is synchronized: on each \texttt{step} the wrapper publishes the
agent's control signals, requests advancement of the federated clock, and blocks
until the federation reaches the next tick before reading the resulting
observations. Third, reward functions are pluggable callables selected by name from
the configuration, so the optimization objective can be changed without touching
code. Crucially, the user never writes the agent's HELICS wiring; the compiler
synthesizes it from the RL configuration block.

\subsection{Standardized telemetry and monitoring}
Because execution is distributed, observability is centralized. Federates buffer
their time series and persist them at the end of a run, partitioned by training and
test phases, and RL federates additionally record the observation seen before and
after each action together with per-episode rewards and lengths. Results can be
streamed to a time-series database and inspected through a ready-made dashboard
that visualizes both physical quantities (for example, indoor temperature against
its comfort band) and learning metrics (for example, episode return), independently
of the application domain.

%==============================================================================
\section{Episodic RL over non-resettable co-simulations}\label{sec:reset}
%==============================================================================
Gymnasium assumes that \texttt{reset()} returns the environment to an initial
state cheaply and arbitrarily often. A co-simulation violates this assumption: each
federate holds internal state, some FMUs do not support being rewound at all, and
tearing down and restarting a federation between episodes is expensive and
disruptive to the shared clock. Episodic RL must therefore be reconciled with a
simulation whose clock only advances. We treat this explicitly as the
\emph{reset problem for non-resettable co-simulations} and expose it as a
first-class, configurable mechanism rather than hiding it.

CosimGym offers four reset strategies, summarized in Table~\ref{tab:reset}. A
\emph{full} reset restores models to their initial conditions and is the closest
analogue to a classical environment reset, at the cost of discontinuity in the
co-simulation. A \emph{soft} reset keeps the current physical state and simply
begins a new episode boundary, which is appropriate when the process is naturally
continuing. A \emph{rolling} reset advances the starting point forward by a fixed
window each episode, sweeping the agent across the simulated horizon and improving
coverage of operating conditions. A \emph{random} reset draws starting conditions
from configured ranges, supporting domain randomization and robustness studies.

\begin{table}[t]\centering
\caption{Episode-reset strategies for non-resettable co-simulations.}
\label{tab:reset}
\small
\begin{tabular}{@{}lll@{}}
\toprule
Strategy & State handling & Typical use \\
\midrule
full    & restore initial conditions & classical episodic training \\
soft    & keep current state         & naturally continuing processes \\
rolling & shift start by a window     & broad coverage of the horizon \\
random  & sample within bounds        & robustness / domain randomization \\
\bottomrule
\end{tabular}
\end{table}

At every episode boundary the framework also handles temporal causality
consistently: staged \texttt{next\_step} inputs from the previous episode are
cleared, current values are re-read, and configured reset defaults are applied so
that the first action of a new episode is chosen from the reset state rather than
from the terminal state of the previous one. Beyond plumbing, this mechanism is a
reusable methodological contribution: the same problem arises for any streaming or
irreversible simulator, including real-time digital twins and physical plants, so
the strategies studied here transfer to sim-to-real settings.

%==============================================================================
\section{Usability and illustrative validation}\label{sec:examples}
%==============================================================================
This section illustrates correctness, breadth, and ease of use; exhaustive control
results are the subject of a companion application study. Unless noted, each
example corresponds to a scenario shipped with the framework, and every figure is
regenerated by the reproducibility harness from its scenario file and a fixed seed.
% TODO: confirm harness + fill all numbers/figures after running experiments.

\subsection{Rapid iteration and swapability}
Replacing a rule-based or PID baseline with an RL agent, or changing a physical
parameter, is a change to the YAML scenario rather than to code; the framework
automatically remaps action bounds and discrete spaces and translates agent
outputs back into engineering units. Listing~\ref{lst:diff} shows the essential
difference between a baseline co-simulation and its RL counterpart.

\begin{lstlisting}[language=,caption={Baseline-to-RL change is declarative (illustrative).},label=lst:diff,basicstyle=\ttfamily\footnotesize,frame=single]
- controller: {model_name: simple_pid_controller}
+ reinforcement_learning_config:
+   agent: {model_name: rl_simple_SACsb3, env: {actions: [..modulation]}}
\end{lstlisting}

\subsection{Verification: spring--mass--damper}
A spring--mass--damper system, run as a single federation and again split across
federations, validates orchestration, multi-rate stepping, and the hierarchical
broker. We compare the simulated trajectory against the analytical solution.
% TODO: Figure 3 (simulated vs analytical) + max trajectory error; confirm no time drift.

\subsection{RL control on composable physics}
A building zone coupled to a heat pump and a weather feed is controlled first by a
PID baseline and then by DQN and SAC agents, using the same physical models.
% TODO: Figure 4 (temperature + power) + comfort/energy metrics vs PID.

\subsection{Multi-formalism: an EnergyPlus FMU}
The identical RL pipeline drives a building modeled in EnergyPlus and exported as
an FMI~2.0 FMU, with no change to the agent or training code, demonstrating the
catalog abstraction in practice.
% TODO: one figure or metrics row for the FMU-based control.

\subsection{Reset-strategy micro-benchmark}
Using the building task, we compare learning under the four reset strategies of
Section~\ref{sec:reset}. This experiment is, to our knowledge, not expressible in
existing co-simulation RL frameworks.
% TODO: Figure 5 (learning curves full/soft/rolling/random) + sample-efficiency table.

\subsection{Composability}
Finally, we quantify engineering effort: the lines of declarative configuration
and of Python needed to go from a base co-simulation to a trained agent, compared
against the estimated effort of a hand-built HELICS-to-Gym wrapper.
% TODO: Table 3 with the actual counts.

%==============================================================================
\section{Impact}\label{sec:impact}
%==============================================================================
CosimGym lowers the barrier to studying control of cyber-physical energy systems
under realistic physics. By making the co-simulation a declarative artifact, it
enables reproducible experiments that can be shared and re-run, and by abstracting
the model layer it lets researchers reuse and combine components across domains and
formalisms instead of rebuilding simulation infrastructure for each study. The
catalog and scenario compiler are reusable beyond the examples shown here: any
model that conforms to the minimal interface, or any FMU, becomes a composable
environment building block. The framework is also suitable for teaching, since a
complete distributed co-simulation can be specified and run from a single readable
file. More broadly, it provides a substrate on which several open research
directions---multi-agent and multi-timescale learning, sim-to-real transfer, and
the systematic generation of diverse training environments---can be pursued without
re-engineering the foundations (Section~\ref{sec:future}).

%==============================================================================
\section{Limitations, trade-offs and future potential}\label{sec:future}
%==============================================================================
\subsection{Limitations and trade-offs}
The generality of the framework has costs. Routing every exchanged variable through
a state store and the HELICS brokers introduces overhead relative to tightly
coupled, single-process simulators, and the appropriate choice of time-step
resolution trades fidelity against communication cost.
% TODO: quantify overhead vs a pure-Python baseline if possible.
Operationally, the framework relies on supporting services (a Redis store and,
optionally, a time-series database) and on HELICS binaries, which adds setup
complexity. The current implementation focuses on single-agent learning and
single-host execution; these are engineering boundaries rather than architectural
ones, as discussed next.

\subsection{Future potential}
The architecture is deliberately positioned to support developments that are
difficult in existing tools:
\begin{enumerate}
  \item \textbf{Multi-agent and federated learning}, with one RL federate per
        subsystem and inter-agent communication modeled explicitly over HELICS,
        including latency and partial observability;
  \item \textbf{multi-timescale and hierarchical RL} that exploits the native
        multi-rate federation, for example a slow planning agent driving a fast
        tracking agent;
  \item \textbf{sim-to-real transfer}, replacing a simulated federate with a
        real-time hardware federate through a configuration change while keeping
        the agent and wiring unchanged;
  \item \textbf{generalist control agents} trained by domain randomization over the
        model catalog, turning the compiler into an environment-generation engine;
  \item \textbf{automatic safety constraints} derived from the typed catalog bounds;
        and
  \item \textbf{automated FMU introspection} and distributed, cluster-native
        scaling.
\end{enumerate}

%==============================================================================
\section{Conclusion}\label{sec:conclusion}
%==============================================================================
CosimGym offers a declarative, model-agnostic substrate that makes realistic,
multi-domain, multi-rate physics the easy part of a reinforcement-learning study,
so that research effort can be directed at the agent rather than at simulation
plumbing. By unifying co-simulation orchestration, model exchange, temporal
synchronization, and episodic interaction behind a single Gymnasium-compatible
interface, and by treating episodic reset over non-resettable simulations as a
first-class concern, the framework provides a reproducible foundation for control
research in cyber-physical energy systems and a basis for the multi-agent,
multi-timescale, and sim-to-real extensions outlined above.

%==============================================================================
% \section*{Acknowledgements} Politecnico di Torino, EC-Lab.
% \appendix \section{Reproducibility} scenario files, seeds, environment, hashes.
\bibliographystyle{elsarticle-num}
% \bibliography{cosimgym}
\end{document}
